\section{The Donninger et al.\ method for ODE-type blowup}
\label{fw:sec}

This section fixes notation and isolates the part of the argument that is
common to all models considered below. The method goes back to Donninger's
work on wave maps \cite{Donninger2011} and was developed into a general
scheme by Donninger and Schörkhuber
\cite{DonningerSchoerkhuber2012,DonningerSchoerkhuber2014,DonningerSchoerkhuber2016,DonningerSchoerkhuber2017},
with later refinements for mode stability \cite{CostinDonningerGlogic2017},
high Sobolev spaces in all dimensions \cite{Ostermann2024,Glogic2025} and
hyperboloidal similarity coordinates \cite{BiernatDonningerSchoerkhuber2021}; see \cite{Donninger2024} for a survey.
We use it here as a black box with five steps, and we point out the (few)
places where the industrial models need something extra.

\subsection{Setting and self-similar variables}
We consider
\begin{equation}\label{fw:eq}
  \partial_t^2 u-\Delta u = F(u) \qquad\text{in } [0,T)\times\mathbb R^d,
\end{equation}
possibly with lower-order terms $\gamma\partial_t u+m^2u$ that will be treated
as perturbations (\cref{fw:subsec-lower}). By finite speed of propagation it
suffices to study \eqref{fw:eq} in the backward light cone
$\Gamma_T=\{(t,x):0\le t<T,\ |x|\le T-t\}$ of the blowup point $(T,0)$; in
particular boundary conditions of the engineering problem (clamped edges of a
MEMS plate, Dirichlet temperature, \dots) are irrelevant as long as the cone
base lies inside the physical domain.

\emph{Similarity coordinates} are
\[
  \tau=-\log(T-t)+\log T,\qquad \xi=\frac{x}{T-t},\qquad (\tau,\xi)\in[0,\infty)\times\overline{B^d},
\]
so that $\partial_\tau=(T-t)\partial_t-x\cdot\nabla_x$ and
$\nabla_\xi=(T-t)\nabla_x$. We write $\Lambda=\xi\cdot\nabla_\xi$.
Suppose the nonlinearity has a scaling exponent $s\in\mathbb R$, i.e.\
$u\mapsto \mu^{s}u(\mu t,\mu x)$ maps solutions to solutions. This holds for
$F(u)=|u|^{p-1}u$ with $s=2/(p-1)$ and for the MEMS force $F(w)=-\lambda
w^{-q}$ with $s=-2/(q+1)$ (a \emph{negative} exponent: the gap $w$ goes to
zero). Setting
\begin{equation}\label{fw:psi}
  \psi_1(\tau,\xi)=(T-t)^{s}\,u(t,x),\qquad \psi_2(\tau,\xi)=(T-t)^{s+1}\,\partial_tu(t,x),
\end{equation}
equation \eqref{fw:eq} becomes the autonomous first-order system
\begin{equation}\label{fw:system}
  \partial_\tau\begin{pmatrix}\psi_1\\ \psi_2\end{pmatrix}
  =\underbrace{\begin{pmatrix}-\Lambda\psi_1-s\psi_1+\psi_2\\ \Delta\psi_1-\Lambda\psi_2-(s+1)\psi_2\end{pmatrix}}_{=:\ \mathbf L_0^{(s)}\Psi}
  +\begin{pmatrix}0\\ F(\psi_1)\end{pmatrix}.
\end{equation}
For the exponential nonlinearity $F(u)=\delta e^u$ one instead sets
$\psi_1=u+2\log(T-t)$, $\psi_2=(T-t)\partial_tu$; the linear part is then
$\mathbf L_0^{(0)}$ and an inhomogeneity $-2$ appears in the first component.
The elementary but useful identity
\begin{equation}\label{fw:shift}
  \mathbf L_0^{(s)}=\mathbf L_0^{(0)}-s\,\mathbf I
\end{equation}
shows that the scaling exponent only \emph{shifts} the free spectrum; this is
what makes the MEMS case ($s<0$) structurally identical to the classical one.

\subsection{ODE blowup and a universal spectrum}
\label{fw:subsec-ode}
An \emph{ODE blowup} solution is a spatially homogeneous solution
$u_T(t)=\kappa(T-t)^{-s}$ (respectively $u_T=\log\kappa-2\log(T-t)$), i.e.\ a
constant fixed point $\Psi_*=(\kappa,s\kappa)$ of \eqref{fw:system}. Its
symmetry orbit consists of
\begin{itemize}[nosep]
\item time translations $T\mapsto T'$ (in the fixed coordinates $(\tau,\xi)$
  adapted to $T$, this is a genuine motion of the fixed point);
\item Lorentz boosts: for $a\in\mathbb R^d$, $|a|<1$,
  \begin{equation}\label{fw:boost}
  u_{T,a}(t,x)=\kappa\,(1-|a|^2)^{s/2}\,\big(T-t+a\cdot x\big)^{-s},\qquad
  \Psi_a(\xi)=\kappa(1-|a|^2)^{s/2}(1+a\cdot\xi)^{-s},
  \end{equation}
  which is again an exact self-similar solution, smooth on
  $\overline{B^d}$ (direct substitution: with $f=T-t+a\cdot x$ one has
  $\partial_t^2f^{-s}-\Delta f^{-s}=s(s+1)(1-|a|^2)f^{-s-2}$);
\item space translations of the blowup point, which do not change $u_T$.
\end{itemize}
Linearising $F(\psi_1)$ at $\kappa$ gives the constant potential
$V_0=F'(u_T)(T-t)^2$. The operator $\mathbf L=\mathbf
L_0^{(s)}+\mathbf L'$, $\mathbf L'\Psi=(0,V_0\psi_1)$, has an eigenvalue
$\lambda$ with eigenfunction $(g,(\lambda+s+\Lambda)g)$ iff
$v=(T-t)^{-s-\lambda}g(x/(T-t))$ solves $\partial_t^2v-\Delta v=V_0(T-t)^{-2}v$,
i.e.\ with $a:=s+\lambda$,
\begin{equation}\label{fw:eig}
  (\delta_{ij}-\xi_i\xi_j)\partial_i\partial_jg-2(a+1)\Lambda g-\big(a(a+1)-V_0\big)g=0\qquad\text{in }B^d.
\end{equation}

\begin{lemma}[$V_0$ is forced by the time-translation mode]\label{fw:lem-V0}
For every ODE blowup with scaling exponent $s$ one has $V_0=(s+1)(s+2)$.
\end{lemma}
\begin{proof}
$\partial_Tu_T\propto (T-t)^{-s-1}$ solves the linearised equation, so
$a=s+1$ (i.e.\ $\lambda=1$, $g\equiv1$) solves \eqref{fw:eig}:
$(s+1)(s+2)-V_0=0$.
\end{proof}
Thus $a^2+a-V_0=(a-a_1)(a-a_2)$ with $a_1=s+1$, $a_2=-(s+2)$, and for
$F=|u|^{p-1}u$ one recovers $V_0=p\kappa^{p-1}=ps(s+1)$, for MEMS
$V_0=q\lambda c^{-q-1}=2q(q-1)/(q+1)^2$ and for $F=\delta e^u$ $V_0=2$.

Writing $g(\xi)=|\xi|^\ell Y_\ell(\xi/|\xi|)\,h(|\xi|^2)$ with a spherical
harmonic $Y_\ell$, \eqref{fw:eig} becomes Gauss' hypergeometric equation
\begin{equation}\label{fw:hyp}
  z(1-z)h''+\big(c-(\alpha+\beta+1)z\big)h'-\alpha\beta h=0,\quad
  c=\tfrac d2+\ell,\quad \alpha=\tfrac{a+\ell-a_1}{2},\quad \beta=\tfrac{a+\ell-a_2}{2},
\end{equation}
(verified symbolically in \texttt{code/}), with Frobenius exponents
$\{0,1-c\}$ at $z=0$ and $\{0,\,c-\alpha-\beta\}$ at $z=1$, where
$c-\alpha-\beta=\frac{d-1}{2}-s-\lambda$.

\begin{theorem}[Spectrum of ODE blowup]\label{fw:thm-spec}
Let $k\in\mathbb N$, $k>\frac d2-s$, and set $\omega_k:=k-\frac d2+s>0$. A
number $\lambda$ with $\operatorname{Re}\lambda>-\omega_k$ is an eigenvalue
of $\mathbf L$ on $\mathcal H^k:=H^k\times H^{k-1}(B^d)$ if and only if
\[
  \lambda\in\{\,1-\ell-2n\,\}\ \cup\ \{\,-2s-2-\ell-2n\,\},\qquad \ell,n\in\mathbb N_0,
\]
and the eigenfunctions are $|\xi|^\ell Y_\ell\cdot(\text{polynomial in }|\xi|^2)$.
In particular, for $s>-1$ the only eigenvalues with $\operatorname{Re}\lambda\ge0$
are $\lambda=1$ ($\ell=0$, time translation, $g\equiv1$) and $\lambda=0$
($\ell=1$, $d$-fold, Lorentz boosts, $g=\xi_j$), and the spectral gap is
\begin{equation}\label{fw:gap}
  \omega_0(s)=\min\{1,\;2s+2\}.
\end{equation}
\end{theorem}
\begin{proof}[Proof sketch]
Since $\Gamma(c)\neq0$ and the second Frobenius solution at $z=0$ behaves like
$|\xi|^{2-d-2\ell}$ (or $\log|\xi|$, or $|\xi|$ when $d=1,\ell=0$), an
$H^k$ eigenfunction ($k\ge2$) must be a multiple of
$h={}_2F_1(\alpha,\beta;c;z)$. At $z=1$ the second exponent produces
$(1-|\xi|^2)^{\gamma}$, $\gamma=\frac{d-1}{2}-s-\lambda$, possibly with a
logarithm; such a function lies in $H^k(B^d)$ only if
$\operatorname{Re}\gamma>k-\frac12$ or $\gamma\in\mathbb N_0$ without log term.
For $\operatorname{Re}\lambda>-\omega_k$ we have $\operatorname{Re}\gamma<k-\frac12$,
so $h$ must be analytic at $z=1$. By the $z\mapsto1-z$ connection formulae \cite[\S15.8]{DLMF} (including the
degenerate cases $\gamma\in\mathbb Z$, in which the coefficients of the logarithm
and of $(1-z)^{\gamma}$, $\gamma<0$, are proportional to
$1/(\Gamma(\alpha)\Gamma(\beta))$ resp.\ $1/(\Gamma(\alpha-m)\Gamma(\beta-m))$)
this happens iff $\alpha\in-\mathbb N_0$ or $\beta\in-\mathbb N_0$, in which case $h$ is a
polynomial. Solving $\alpha=-n$, $\beta=-n$ for $\lambda=a-s$ gives the two
families. For $s>-1$ the second family is $\le -2s-2<0$, and $1-\ell-2n\ge0$
only for $(\ell,n)\in\{(0,0),(1,0)\}$; the next eigenvalues are $-1$
($\ell=0,n=1$ and $\ell=2,n=0$) and $-2s-2$, whence \eqref{fw:gap}.
\end{proof}

\begin{remark}
\Cref{fw:thm-spec} contains, as special cases, the spectral statements of
\cite{DonningerSchoerkhuber2012,DonningerSchoerkhuber2014,DonningerSchoerkhuber2017,Ostermann2024}
for $F=|u|^{p-1}u$, $p>1$ ($s=2/(p-1)$, gap $\min\{1,\frac{2(p+1)}{p-1}\}=1$),
and it extends verbatim to the MEMS range $p=-q<-1$
($s=-\frac{2}{q+1}\in(-1,0)$, gap $\min\{1,\frac{2(q-1)}{q+1}\}$) and to the
exponential case ($s=0$, gap $1$). The constraint $s>-1$ is exactly the
condition $q>1$ for MEMS touchdown to be of self-similar ODE type
(\cref{mems:sec}). Semisimplicity of $\lambda=1$ and $\lambda=0$ (no Jordan
chains) is checked by showing that $(\mathbf L-\lambda)\mathbf u=\mathbf
g_\lambda$ has no $\mathcal H^k$ solution, an ODE computation of the same
type as in \cite{DonningerSchoerkhuber2017,Ostermann2024}; it is carried
out model by model below.
\end{remark}

\subsection{The five steps}
\label{fw:subsec-steps}
\paragraph{Step 1: free evolution.} By \eqref{fw:shift} it suffices to
understand $\mathbf L_0^{(0)}$. Heuristically, for a free wave
$\|\nabla_\xi^k\psi_1(\tau)\|_{L^2(B^d)}=(T-t)^{k-d/2+s}\|\nabla_x^ku(t)\|_{L^2(B_{T-t})}$,
and the last factor is bounded by local energy conservation, which suggests
the decay rate $\omega_k=k-\frac d2+s$ on top-order seminorms. The low-order
part is not dissipative: polynomial free waves of degree $m$ give eigenvalues
$-s-m$ of $\mathbf L_0^{(s)}$, and for the MEMS case $s<0$ the constant
wave $w\equiv1$ gives the \emph{positive} eigenvalue $-s=\sigma$. The rigorous
version (inner products on $\mathcal H^k$ built from $\dot H^k$-type
seminorms plus boundary/low-order terms, Lumer--Phillips) is
\cite{Ostermann2024} and \cite{Glogic2025}: there is an equivalent
inner product and a finite-rank operator $\mathbf K_0$ such that
$\mathbf L_0^{(s)}-\mathbf K_0$ generates a semigroup of growth bound
$\le-\omega$ for any prescribed $\omega<\omega_k$ (after possibly increasing
$k$). Since the finite-rank part is absorbed into the compact perturbation in
Step~2, the positive free eigenvalue in the MEMS case is harmless.

\paragraph{Step 2: compact perturbation.} $\mathbf L'+\mathbf K_0$ is compact
on $\mathcal H^k$ (Rellich), hence $\mathbf L$ generates a semigroup
$\mathbf S(\tau)$ with the same essential growth bound $\le-\omega_k$, and
$\sigma(\mathbf L)\cap\{\operatorname{Re}\lambda>-\omega_k\}$ consists of
isolated eigenvalues of finite algebraic multiplicity
\cite[Ch.~IV]{EngelNagel2000}.

\paragraph{Step 3: mode stability.} Determine all eigenvalues in
$\operatorname{Re}\lambda>-\omega_0$; for the ODE profile this is
\cref{fw:thm-spec}. For non-explicit profiles (wave maps, Yang--Mills,
non-trivial profiles of supercritical equations) this is the hard step and
is done by supersymmetric removal of symmetry modes and the quasi-solution
method of Costin, Donninger and Glogić \cite{CostinDonningerGlogic2017}.

\paragraph{Step 4: linear stability.} With the Riesz projection $\mathbf P$
onto the symmetry eigenvalues ($\lambda=1$, and $\lambda=0$ in the nonradial
case), $\|\mathbf S(\tau)(\mathbf I-\mathbf P)\|\le M_\omega e^{-\omega\tau}$ for
every $\omega<\min\{\omega_0,\omega_k\}$, because the growth bound of a
semigroup equals the maximum of its spectral bound and its essential growth
bound.

\paragraph{Step 5: nonlinear stability.} Writing $\Psi=\Psi_*+\Phi$, the
perturbation solves $\partial_\tau\Phi=\mathbf L\Phi+\mathbf N(\Phi)$ with
$\mathbf N(\Phi)=(0,F(\kappa+\phi_1)-F(\kappa)-F'(\kappa)\phi_1)$. If $k>d/2$,
$H^k(B^d)$ is an algebra embedded in $L^\infty$, so for smooth $F$ near
$\kappa$ (for MEMS: away from $w=0$) one has
$\|\mathbf N(\Phi)-\mathbf N(\tilde\Phi)\|\lesssim(\|\Phi\|+\|\tilde\Phi\|)\|\Phi-\tilde\Phi\|$.
The Duhamel equation is solved in
$\mathcal X_\omega=\{\Phi:\sup_\tau e^{\omega\tau}\|\Phi(\tau)\|_{\mathcal H^k}<\infty\}$
after subtracting the Lyapunov--Perron correction
$\mathbf C(\Phi,\mathbf u)=\mathbf P\mathbf u+\int_0^\infty e^{-\tau'}\mathbf P_1\mathbf N(\Phi(\tau'))\,d\tau'
+\int_0^\infty\mathbf P_0\mathbf N(\Phi(\tau'))\,d\tau'$. The correction is then
removed by choosing the blowup time $T^*$ (and, nonradially, the boost
parameter $a^*$ and the blowup point) via the Brouwer fixed point theorem:
the derivative of the rescaled data with respect to $(T,a)$ at the
background is a non-degenerate map onto $\operatorname{ran}\mathbf P$, because
the symmetry eigenfunctions are precisely the generators of these
transformations. The outcome is a statement of the following form.

\begin{theorem}[template]\label{fw:thm-template}
Let $k>\frac d2+\max\{0,-s\}$ be large enough and $0<\omega<\omega_0$.
There is $\varepsilon>0$ such that for data $(u_0,u_1)$ with
$\|(u_0,u_1)-(u_{1,0}(0),\partial_tu_{1,0}(0))\|_{H^k\times H^{k-1}(B_{1+\varepsilon})}\le\varepsilon$
there exist $T^*\in(1-\varepsilon,1+\varepsilon)$ and
$a^*\in B_\varepsilon$ (with $a^*=0$ for radial data) such that the
solution blows up at $(T^*,0)$ and
\[
  \big\|(T^*-t)^{s}u(t,(T^*-t)\,\cdot)-\Psi_{a^*}\big\|_{H^k(B^d)}
  +\big\|(T^*-t)^{s+1}\partial_tu(t,(T^*-t)\,\cdot)-(\Lambda+s)\Psi_{a^*}\big\|_{H^{k-1}(B^d)}
  \lesssim (T^*-t)^{\omega}.
\]
\end{theorem}
No translation parameter is needed: the ODE profile is translation invariant,
and a space translation of a boosted profile is a time translation of it
($u_{T,a}(t,x-x_0)=u_{T-a\cdot x_0,a}(t,x)$). Applying the theorem at nearby
base points shows that ODE blowup occurs along a spacelike hypersurface
$t=T^*(x)$, with $a^*$ playing the role of its gradient (cf.\ \cref{mems:rem-main}).

\subsection{Lower-order terms}\label{fw:subsec-lower}
Engineering models almost always contain damping or restoring forces. With
$T-t=Te^{-\tau}$, a damping term $\gamma\partial_tu$ adds
$-\gamma Te^{-\tau}\psi_2$ to the second line of \eqref{fw:system}, and a
linear restoring force $m^2u$ adds $-m^2T^2e^{-2\tau}\psi_1$. Two
consequences: (i) the ODE blowup solution is no longer exactly self-similar,
$u(t)=\kappa(T-t)^{-s}(1+O(T-t))$, so the background in similarity variables is
$\Psi_*+O(e^{-\tau})$; (ii) the linearised flow is an \emph{integrable}
($\int_0^\infty\|e^{-\tau}\mathbf B\|\,d\tau<\infty$) bounded perturbation of
$\mathbf S(\tau)$. Both are compatible with Steps~4--5 as long as the decay rate
of the forcing ($1$ for damping) is not below the target rate $\omega$, which
holds since $\omega<\omega_0\le1$. In the radial case only $T$ must be
modulated. In the nonradial case, the mass term preserves Lorentz
invariance, whereas damping breaks it. Even then the neutral $\ell=1$ modes
are tangent to an exact family: $w=W(T-t+a\cdot x)$ solves the damped
equation if $W$ solves the spatially homogeneous equation with damping
$\gamma/\sqrt{1-|a|^2}$ (after rescaling the variable by $\sqrt{1-|a|^2}$).
Moreover, since there is no Jordan block at $\lambda=0$, the neutral
component obeys $\partial_\tau\mathbf P_0\Phi=\mathbf P_0\mathbf G(\Phi)$
with a forcing that is \emph{integrable} in $\tau$; hence $\mathbf P_0\Phi$
has a limit and a \emph{static} choice of the boost parameter in the
Lyapunov--Perron step suffices (\cref{mems:prop-damped}). A dynamical
modulation is only needed when the neutral modes are forced
non-integrably. The details are in \cref{power:sec} and in the model
sections.

\subsection{What the method gives, and what it does not}
The output is \emph{local}: stability of the blowup profile inside the
backward light cone of the (moving) blowup point, for perturbations small in
a high Sobolev norm, with an explicit exponential rate in $\tau$, i.e.\ a
power rate in $T^*-t$. It does not, by itself, determine the global blowup
set or describe the solution outside the cone; for the latter see the
hyperboloidal variant \cite{BiernatDonningerSchoerkhuber2021} and
\cite{ChenEtAl2024}. For the industrial models
we treat, it nevertheless yields exactly the information an engineer would
ask for: the \emph{touchdown/runaway rate}, the \emph{local profile}, and the
\emph{robustness} of both under perturbations of data and of lower-order
model terms.
